\section{Lean formalization architecture against the current API}\label{sec:lean}

\subsection{Target location and imports}

The implementation should be added under
\[
\RepoPath{Logic/PrimitiveRecursiveArithmetic/Derivations/Lean/IPRADerivations/Framework/}
\]
and imported only after the existing core.  It should not copy \texttt{ConciseIPRA} syntax or define a parallel proof calculus.  The proposed root module is \texttt{IPRADerivations/Framework.lean}.

The minimum current imports are:
\begin{itemize}
  \item \texttt{IPRADerivations.Core.Basic} for \texttt{PR}, \texttt{PTm}, \texttt{PRAAx}, \texttt{Prf}, \texttt{IPRA.ind}, and \texttt{IPRA.inst};
  \item \texttt{IPRADerivations.Core.Order} and \texttt{Core.Monus} for arithmetic comparisons;
  \item \texttt{IPRADerivations.Core.Pairing} for \texttt{eqb}, \texttt{pairT}, \texttt{fstT}, \texttt{sndT}, and two-track simultaneous recursion;
  \item \texttt{IPRADerivations.Core.Search} for checked least bounded search;
  \item \texttt{Core.BoundedSums}, \texttt{Core.SumLemmas}, and \texttt{Core.SumReflect} for finite folds;
  \item \texttt{Core.Tactic}, \texttt{Core.Integers}, and \texttt{Core.LinearAlgebra} for algebraic local rules.
\end{itemize}
Application modules import additional checked mathematics only as needed.

\subsection{Proposed module graph}

\begin{enumerate}[label=\textbf{F\arabic*.}]
  \item \texttt{Framework/Kernel.lean}: local rules, proof schedules, compiled conditional theorems, loop certificates.
  \item \texttt{Framework/Bool.lean}: Boolean range, truth tables, reflected atoms and connectives.
  \item \texttt{Framework/BoundedFormula.lean}: generated \texttt{allB}/\texttt{anyB} folds and syntax-directed fixed-formula reflection.
  \item \texttt{Framework/SearchAdapter.lean}: thin success-at-one wrapper around checked \texttt{Search.mu}; bounded argmin and first-valid corollaries.
  \item \texttt{Framework/FiniteCode.lean}: fixed tuples, sequences, finite tables, lookup/update, and proved equations, built on checked pairing.
  \item \texttt{Framework/Trace.lean}: packed state, freeze-after-done, rank-bounded loop certificates.
  \item \texttt{Framework/Residual.lean}: residual normalization metadata and proof-schedule closure.
  \item \texttt{Framework/Demand.lean}: generated acyclic reverse demand and capped cyclic closure certificates.
  \item \texttt{Framework/Relation.lean}: viability table, bounded path checker, least viable path.
  \item \texttt{Framework/Observer.lean}: exact, diameter, and gap-guarded observer elimination.
  \item \texttt{Framework/Rational.lean}: signed rational codes, dyadic error predicates, triangle inequalities, clipping, and finite error sums.
  \item \texttt{Framework/Compiler.lean}: metalevel graph records, query specialization, source generation, and provenance.
  \item \texttt{Framework/Benchmarks/*.lean}: small closed examples, then rational cubic bisection.
\end{enumerate}
The archive's Lean files are source-oriented blueprints for these modules.  They are not claimed to compile in this environment.

\subsection{Kernel-facing core}

The principal abstraction should remain small enough that its soundness is transparent.

\begin{lstlisting}
namespace LeanProofs.IPRA.Framework

open LeanProofs.ConciseIPRA
open Syn
open scoped Fm

structure Rule (n : Nat) where
  premises : List (Fm PTm n)
  conclusion : Fm PTm n
  sound : premises |- [PRAAx] conclusion

namespace Rule

theorem use {Gamma : List (Fm PTm n)} (r : Rule n)
    (supply : forall psi, psi in r.premises ->
      Gamma |- [PRAAx] psi) :
    Gamma |- [PRAAx] r.conclusion :=
  r.sound.trans supply

end Rule

inductive Schedule (n : Nat) (Gamma : List (Fm PTm n)) :
    Fm PTm n -> Type
  | done (h : Gamma |- [PRAAx] phi) : Schedule n Gamma phi
  | have_ (hpsi : Gamma |- [PRAAx] psi)
      (rest : Schedule n (psi :: Gamma) phi) :
      Schedule n Gamma phi

noncomputable def Schedule.compile :
    Schedule n Gamma phi -> Gamma |- [PRAAx] phi
  | .done h => h
  | .have_ hpsi rest =>
      (Schedule.compile rest).trans
        (List.forall_mem_cons.2
          <| And.intro hpsi (fun chi hchi => Prf.hyp hchi))

end LeanProofs.IPRA.Framework
\end{lstlisting}
The exact proof above matches the pattern already used successfully in the earlier source outline.  It should be the first compilation gate.

A residualized endpoint is simply a proof-bearing record, not an axiom mechanism:

\begin{lstlisting}
structure ConditionalTheorem (n : Nat) where
  residuals : List (Fm PTm n)
  permanent : List (Fm PTm n)
  conclusion : Fm PTm n
  proof : residuals ++ permanent |- [PRAAx] conclusion
\end{lstlisting}
A closure theorem takes derivations of every residual formula and applies \texttt{Prf.trans}.  The metalevel compiler may store names, keys, and provenance separately; only the formula list and proof enter the trusted endpoint.

\subsection{Loop certificates}

The checked \texttt{IPRA.ind} interface should be used directly.

\begin{lstlisting}
structure LoopCertificate (k : Nat) where
  inv : Fm PTm (k + 1)
  qf : inv.QF
  zeroCase : ([] : List (Fm PTm k)) |- [PRAAx]
    inv.subst (Syn.inst PTm.zero)
  succCase : [inv] |- [PRAAx]
    inv.subst (Syn.repl (PTm.succ (Syn.var 0)))

namespace LoopCertificate

def compile (c : LoopCertificate k) :
    ([] : List (Fm PTm (k + 1))) |- [PRAAx] c.inv :=
  LeanProofs.IPRA.ind c.inv c.qf c.zeroCase c.succCase

theorem at (c : LoopCertificate k)
    (sigma : Fin (k + 1) -> PTm n) :
    Gamma |- [PRAAx] c.inv.subst sigma :=
  LeanProofs.IPRA.inst c.compile sigma

end LoopCertificate
\end{lstlisting}
The main engineering rule is that the generated invariant contains only equations, order formulas, and Boolean equations.  Bounded quantifiers are represented by terms before this structure is built.

\subsection{Boolean reflection}

A fixed reflected formula can be represented by:

\begin{lstlisting}
structure Reflected (n : Nat) where
  formula : Fm PTm n
  code : PTm n
  range : ([] : List (Fm PTm n)) |- [PRAAx]
    (code = PTm.zero) or (code = PTm.one)
  iffOne : ([] : List (Fm PTm n)) |- [PRAAx]
    Fm.iff formula (code = PTm.one)
\end{lstlisting}
Here \texttt{PTm.one} is a local abbreviation for \texttt{succ zero}.  Equality reflection should reuse the checked \texttt{eqb} theorem from \texttt{Core.Pairing}.  Order reflection should be proved against the current \texttt{le}/\texttt{lt} definitions.  Boolean connectives are fixed arithmetic selectors with four-case truth-table derivations.

For a fixed predicate term \(p(i,x)\), generate \texttt{allB} and \texttt{anyB} symbols with \texttt{PR.prec}.  Prove:
\begin{enumerate}[label=(\arabic*)]
  \item zero and successor equations;
  \item Boolean range;
  \item endpoint introduction/elimination lemmas such as
  \[
    \mathsf{allB}(N,x)=1\wedge i<N\to p(i,x)=1;
  \]
  \item witness introduction for \texttt{anyB}; and
  \item full bounded-quantifier wrappers by combining those lemmas with checked search.
\end{enumerate}
The object induction matrices can remain QF because \(i\) is a free parameter and the fold is a term.

The syntax-directed compiler is a Lean metalevel recursion on one concrete \texttt{Fm} value.  It returns a \texttt{Reflected}; it does not define a universal object-level evaluator for formula codes.

\subsection{Search adapter: reuse, do not reprove}

The current checked search API has the schematic form:

\begin{lstlisting}
namespace LeanProofs.IPRA.Search

abbrev mu (f : PTm 2) (N x : PTm n) : PTm n

theorem search_spec (N i x : PTm n) :
  Gamma |- [PRAAx]
    le (mu f N x) N and
      ((lt (mu f N x) N -> inst2 f (mu f N x) x = zero) and
       (lt i (mu f N x) -> not (inst2 f i x = zero)))

theorem search_found
    (hf : Gamma |- [PRAAx] inst2 f i x = zero)
    (hi : Gamma |- [PRAAx] lt i N) :
  Gamma |- [PRAAx]
    lt (mu f N x) N and
      ((inst2 f (mu f N x) x = zero) and le (mu f N x) i)

end LeanProofs.IPRA.Search
\end{lstlisting}
The framework module should define only adapters.  For a Boolean success predicate \(p\), let \(f=1\dotminus p\).  Prove once that under the range theorem, \(f=0\leftrightarrow p=1\).  Then derive \texttt{firstValid}, \texttt{firstValid\_sound}, \texttt{firstValid\_minimal}, bounded argmin, and first covering chart by instantiating the checked theorems.  The old \texttt{SearchBlueprint.lean} should therefore be replaced, not completed as an independent recurrence.

\subsection{Finite codes and packed state}

Use the current checked pairing
\[
  \langle a,b\rangle=(a+b)^2+a
\]
with \texttt{fstT}/\texttt{sndT} projection theorems.  A fixed \(k\)-tuple is a right-nested pair; generate named projections and prove the projection equations by repeated congruence and \texttt{fst\_pair}/\texttt{snd\_pair}.  This is enough for fixed graph states.

Dynamic finite sequences need a length-tagged list code and PR lookup/update.  Implement them only after fixed tuples compile.  Required public theorems are:
\begin{itemize}
  \item lookup at a written index;
  \item lookup at other indices is unchanged;
  \item update preserves length and out-of-range default;
  \item a bounded fold over lookup agrees with its recursive accumulator;
  \item all formulas used in these theorems have QF matrices.
\end{itemize}

For two mutually recursive tracks, reuse checked \texttt{sim\_1}/\texttt{sim\_2}.  For more tracks, either nest the pair state into one recursion or prove a generated fixed-arity wrapper.  No new simultaneous-recursion axiom is needed.

\subsection{Residual compilation}

The metalevel graph record separates data dependencies, proof dependencies, and residual boundaries:

\begin{lstlisting}
structure Boundary (n : Nat) where
  key : String
  formula : Fm PTm n
  origin : String

structure Node (n : Nat) where
  name : String
  dataParents : Array Nat
  proofParents : Array Nat
  boundaries : Array (Boundary n)
  conclusion : Fm PTm n
  build : NodeEnvironment n -> Rule n
\end{lstlisting}
The precise \texttt{build} type may be syntax-generating rather than an ordinary function, but its output must contain a \texttt{Rule.sound}.  Normalize residuals by backward reachability and stable keys.  Generate a \texttt{Schedule} whose base context is \texttt{residuals ++ permanent}.  The theorem \texttt{ConditionalTheorem.close} supplies boundary proofs by \texttt{Prf.trans}.

No theorem about syntactic graph correctness has to be trusted initially: the generated proof term is independently checked.  A verified internal compiler can be a later optimization.

\subsection{Capped demand closure}

Implement two levels.

\paragraph{Semantic prototype.}
Prove in ordinary Lean, for \texttt{Fin N -> Nat}, that a monotone inflationary capped update stabilizes after
\[
  1+\sum_i(C_i-s_i)
\]
steps and yields the least pre-fixed point.  This validates the algorithm and gives executable tests; it is not yet an iPRA proof.

\paragraph{Object certificate.}
For each generated fixed system, emit PR terms for the vector state, update, rank, overflow flag, and bounded iteration.  The generator supplies a \texttt{LoopCertificate} with:
\begin{itemize}
  \item seed/cap preservation;
  \item componentwise inflation;
  \item changed \(=1\) implies rank increases;
  \item rank is at most the cap sum;
  \item once changed \(=0\), the state freezes.
\end{itemize}
At terminal fuel derive fixedness.  Prove leastness by a second QF induction with an arbitrary candidate vector as free parameters.  Export raw edge adequacy under an object proof that overflow is zero.  The cap proof may itself be a residual formula.

\subsection{Relational viability}

The first formal target should have a fixed rectangular width \(B\) and variable horizon \(H\).  Use a coded Boolean table with rows \(0,\ldots,H\).  The reverse recurrence is
\begin{align*}
 W_H(a)&=V_H(a)\wedge T(a),\\
 W_i(a)&=V_i(a)\wedge \mathsf{Any}_{b<B}
       (E_i(a,b)\wedge W_{i+1}(b)).
\end{align*}
Generate a reverse loop writing one row per step.  Its QF invariant says that all written cells satisfy the recurrence and unwritten cells have the default.  Prove a bounded coded-path checker and the equivalence between \(W_i(a)=1\) and existence of an accepted suffix code.  Use the checked search adapter to select the least viable start and successor.

Suggested theorem endpoints are:

\begin{lstlisting}
theorem viable_iff_suffix :
  Gamma |- [PRAAx] Fm.iff
    (viable H i a = one)
    (pathExistsFormula H i a)

theorem selectedPath_valid :
  (anyStartViable H = one :: Gamma) |- [PRAAx]
    pathCheck H (selectedPath H) = one
\end{lstlisting}
The bounded existential formula is optional in internal consumers: the Boolean equation itself is often the cleaner certificate.

\subsection{Stable observers}

Begin with explicit candidate terms rather than binder-heavy records.

\begin{lstlisting}
theorem exactObserver_of_selected
    (hsel : Gamma |- [PRAAx] valid selected)
    (hstable : valid selected :: valid arbitrary :: Gamma
      |- [PRAAx] obs selected = obs arbitrary) :
    valid arbitrary :: Gamma |- [PRAAx]
      obs selected = obs arbitrary
\end{lstlisting}
Then package totality plus checked search to hide \texttt{selected}.  For diameter observers, represent absolute rational inequalities by cross-multiplied natural equations/order formulas.  Prove triangle inequality once, then finite telescoping sums through checked sum lemmas.  The regular-name theorem takes an explicit tail-modulus term; it must not infer one from pointwise shrinking alone.

\subsection{Rational and analytic contracts}

Use a total code \((p,m,d)\) representing \((p-m)/(d+1)\).  Addition and multiplication are polynomial.  Reciprocal returns a default at zero and has a separate apartness guard.  Dyadic error
\[
  |r|\le 2^{-k}
\]
is represented by a cross-multiplied natural inequality, avoiding object-language division.  Required shared theorems are:
\begin{itemize}
  \item rational equality/order reflection;
  \item triangle and product error inequalities;
  \item clipping to a rational interval and nonexpansiveness;
  \item finite sum/product error accumulation;
  \item guard-strengthening and demand monotonicity.
\end{itemize}
These support finite nets, interval methods, series truncations, and the cubic benchmark.

\subsection{Benchmark order}

The formalization should proceed through endpoints of increasing strength.
\begin{enumerate}[label=(\arabic*)]
  \item Compile \texttt{Rule.use}, \texttt{Schedule.compile}, and \texttt{ConditionalTheorem.close}; audit axioms.
  \item Build a Boolean predicate and the checked-search adapter; prove a closed least-witness example.
  \item Generate a three-coordinate capped demand closure and prove fixedness/leastness in \(\Prf\,\PRAAx\).
  \item Build a small layered relation with one dead branch; prove the selected path is valid and an exact observer is choice-independent.
  \item Compile a deterministic finite graph combining an inherited algebra theorem, a guard, checked search, and a loop.
  \item Implement rational bisection for \(y^3+y=x\) on \([0,1]\), first for rational \(x\), then for a represented input with an explicit precision schedule.
\end{enumerate}
Only after steps 1--5 should the generic source generator be generalized.

\subsection{Axiom and build gates}

Every completed module should satisfy all of the following.
\begin{enumerate}[label=(\arabic*)]
  \item Public mathematical endpoints are explicit formulas and actual \(\Prf\,\PRAAx\) values.
  \item There is no \texttt{sorry}, theorem-specific \texttt{axiom}, or opaque semantic replacement for a derivation.
  \item Every induction call supplies an explicit \texttt{Fm.QF} proof.
  \item \texttt{\#print axioms} shows no \texttt{sorryAx}; metatheoretic axioms are recorded separately from object-theory assumptions.
  \item The module builds from the pinned commit or a newly recorded re-audit.
  \item Residual assumptions occur in theorem contexts and provenance, not as unproved fields masquerading as data.
  \item Native transfer is applied only after the ordinary derivation is complete.
  \item Python or semantic Lean tests are labeled as implementation evidence, never as object-level proof.
\end{enumerate}
